\section{Introduction}\label{sec:introduction}

Let $\Omega\subset\R^3$ be a bounded connected domain with smooth boundary.
For real functions $\lambda,\mu\in C^\infty(\overline\Omega)$ satisfying
\begin{equation}\label{eq:energy-positivity}
  \mu>0,\qquad 3\lambda+2\mu>0
  \quad\text{on }\overline\Omega,
\end{equation}
define the strain, stress, and static elasticity operator by
\begin{equation}\label{eq:physical-operator}
 e(u)=\frac{\nabla u+(\nabla u)^T}{2},\qquad
 \sigma_{\lambda,\mu}(u)=\lambda(\diverg u)I+2\mu e(u),\qquad
 L_{\lambda,\mu}u=\diverg\sigma_{\lambda,\mu}(u).
\end{equation}
Here $\mu$ is the shear modulus and $\lambda+2\mu/3$ is the bulk modulus.
The inequalities in \eqref{eq:energy-positivity} give uniformly positive
elastic energy. For each $f\in H^{1/2}(\partial\Omega;\R^3)$ there is a
unique $u_f\in H^1(\Omega;\R^3)$ with $L_{\lambda,\mu}u_f=0$ and trace $f$.
The displacement-to-traction map is defined weakly by
\begin{equation}\label{eq:physical-dn}
 \langle\Lambda_{\lambda,\mu}f,g\rangle
 =\int_\Omega\left[
 \lambda(\diverg u_f)(\diverg v_g)+2\mu e(u_f):e(v_g)
 \right]\dd x,
\end{equation}
where $v_g\in H^1(\Omega;\R^3)$ has trace $g$, and $A:B=\tr(A^TB)$.
The weak equation makes this expression independent of the extension.
For smooth boundary data it is the traction
$\sigma_{\lambda,\mu}(u_f)n$, with $n$ the outward normal.

\begin{theorem}\label{thm:main}
Let $\Omega\subset\R^3$ be any bounded connected domain with
$C^\infty$ boundary. For $j=1,2$, let
$\lambda_j,\mu_j\in C^\infty(\overline\Omega;\R)$ satisfy
$\mu_j>0$ and $3\lambda_j+2\mu_j>0$ on $\overline\Omega$.
If
\[
 \Lambda_{\lambda_1,\mu_1}
 =\Lambda_{\lambda_2,\mu_2}
 \colon H^{1/2}(\partial\Omega;\R^3)
       \longrightarrow H^{-1/2}(\partial\Omega;\R^3),
\]
then $\lambda_1=\lambda_2$ and $\mu_1=\mu_2$ throughout $\Omega$.
\end{theorem}

Thus Theorem~\ref{thm:main} gives a positive resolution of the smooth
three-dimensional isotropic elastic Calder\'on uniqueness problem under
\eqref{eq:energy-positivity}. It uses the full zero-frequency boundary
operator in fixed Euclidean coordinates. The coefficients need not be
analytic or close to constants, and their agreement near the boundary is
not an additional hypothesis.

\paragraph{History and significance.}
The inverse problem asks whether static boundary measurements determine
both spatially varying elastic moduli.
The scalar conductivity problem posed by Calder\'on
\cite{Calderon1980} and the complex geometric optics uniqueness theorem
of Sylvester and Uhlmann \cite{SylvesterUhlmann1987} are important
antecedents, but elasticity couples two moduli through vector
displacements. Nakamura and Uhlmann recovered the elastic boundary
Taylor series \cite{NakamuraUhlmann1995Boundary}. Their original global
claim \cite{NU1994} was corrected in \cite{NU2003}. The erratum identifies
two defects: the prescribed initial-value problem for a planar matrix
transport need not be solvable, and substitution of the corrected complex
geometric optics solutions yields a pseudodifferential equation where a
partial differential equation had been asserted. Its corrected theorem
proves uniqueness when both $\|\nabla\mu_j\|_{C^m}$ are sufficiently
small.

Eskin and Ralston developed higher-order matrix complex geometric optics
expansions using invertible planar transport frames, together with
restricted elasticity uniqueness results
\cite{EskinRalston2002Elasticity,EskinRalston2004}. A direct inspection of
the differential map from their auxiliary vector--scalar fields to
physical displacement \cite[Section~3]{EskinRalston2004} shows that it is
noninjective. The reduction alone therefore does not transfer equality
of the physical boundary maps to equality of unrestricted auxiliary
Cauchy data. The proof here transfers exact physical solutions while
retaining their actual divergence.

In two dimensions, Imanuvilov and Yamamoto proved global uniqueness for
smooth coefficients satisfying $\mu>0$ and $\lambda+\mu>0$, without a
smallness assumption \cite{ImanuvilovYamamoto2015Elasticity}.
In three dimensions, Lin and Nakamura gave local boundary reconstruction
at finite regularity \cite{LinNakamura2017}; Tan and Liu obtained boundary
jets on smooth Riemannian manifolds and global recovery under their
analytic hypotheses
\cite{TanLiu2023}. A recent preprint of Chen, Jiang, Liu, and Tao treats
structured shear moduli under axisymmetry, separation, or directional
quasianalyticity in the specified geometries
\cite{ChenJiangLiuTao2026Elasticity}.
Theorem~\ref{thm:main} treats the full smooth class under
\eqref{eq:energy-positivity}.

\paragraph{Conventions.}
All differential identities are complexified by linearity. Dot products
of complex vectors, including the null condition and orthogonality, use
the complex bilinear extension of the Euclidean product. Sobolev norms
and Hilbert-space adjoints use the usual Hermitian structure.
Equality of the real boundary maps extends complex linearly to the
complex boundary data used in the proof.

\paragraph{Proof strategy.}
The proof converts equality of the boundary maps into an exact exterior
comparison of physical solutions. Boundary determination first gives
matching boundary jets, from which we construct common smooth extensions
of the coefficients outside $\Omega$.
Given a solution in the first medium, equality of displacement and
traction then lets us glue a solution in the second medium to it across
$\partial\Omega$. The two solutions, and their actual divergences,
coincide in the exterior. Section~\ref{sec:boundary-reduction} establishes
this transfer.

Fix a nonzero complex null vector $\theta$, so
$\theta\cdot\theta=0$, and write $D_\theta=\theta\cdot\nabla$.
The real and imaginary parts of $\theta$ span a plane on which
$D_\theta$ is a nonzero multiple of a Cauchy--Riemann operator; the
remaining real coordinate is transverse to these planes.
For solutions with exponential phase $e^{\tau\theta\cdot x}$,
the leading displacement--divergence pair $(a,b)$ has
$a\in H_\theta=\{a\in\C^3:\theta\cdot a=0\}$ and $b\in\C$.
The resulting transport space
$\Ecal_\theta=H_\theta\oplus\C$ has dimension three.
The displacement--divergence reduction to a system with Laplacian
principal part is classical in Lam\'e unique continuation
\cite{Weck1969,ImanuvilovYamamoto2004Carleman}. The formulation here uses
an independent vector--scalar pair $(u,d)$, retains that principal part
after normalization, and satisfies a compatibility identity for the
defect $\diverg u-d$.
Its supported Carleman estimate, together with a compatible high-order
transport recursion and a correction using the physical elasticity
operator, realizes a full local frame in $\Ecal_\theta$ by exact
displacements. The construction controls both the leading pair and its
weighted first derivatives. Sections~\ref{sec:analytic-estimates}
and~\ref{sec:physical-amplitudes} prove these facts.

The leading equations have a useful normalization. For each coefficient
pair, set
\[
 p=\sqrt\mu\,a,\qquad
 s=-\frac{\lambda+\mu}{2\sqrt\mu}\,b
       -\frac12\nabla\log\mu\cdot p.
\]
This is an invertible change of variables on $\Ecal_\theta$, since
$\mu>0$ and $\lambda+\mu>0$. In the $j$th medium, the normalized
transport takes the form
\[
 D_\theta\binom{p}{s}
   =M_{j,\theta}\binom{p}{s},\qquad
 M_{j,\theta}=
 \begin{pmatrix}0&\theta\\Q_{j,\theta}\cdot&T_{j,\theta}\end{pmatrix}.
\]
Here $Q_{j,\theta}\cdot$ is the covector
$p\mapsto Q_{j,\theta}\cdot p$. The lower coefficients are smooth in $x$
and linear in $\theta$;
Lemma~\ref{lem:transport-normal-form} gives their formulas.
The upper equation $D_\theta p=\theta s$ is the same for both media.

Transport comparison followed by holomorphic variation of the complex
direction has a close precedent in Ceki\'c's work on connection Laplacians
\cite{Cekic2025}, whose parameter argument follows Eskin's analysis of
Yang--Mills potentials \cite{Eskin2001}. Ceki\'c treats full
connection-system boundary data on a trivial vector bundle. The comparison
needed here must be obtained from physical elasticity data on the
constrained fibers $\Ecal_\theta$.

On a cylinder given by a planar disk times a transverse interval, let
$Y_1:\C^3\to\Ecal_\theta$ be the matrix of normalized columns of one
local frame realized in the first medium. Transfer its exact physical
solutions to the second
medium. The two augmented operators have the same principal part, so the
forcing for the difference of the physical pairs has only first order.
The derivative bound and the supported Carleman estimate make the
transferred leading pairs bounded in $L^2$. Weak limits supply a matrix
$Y_2:\C^3\to\Ecal_\theta$ of three second-medium transport columns.
The two normalizations agree in
the exterior because the coefficients do, so $Y_2=Y_1$ there.
The local comparison $G=Y_2Y_1^{-1}$ is therefore an endomorphism of
$\Ecal_\theta$ equal to the identity on the part of that cylinder
outside $\overline\Omega$.
Only the original columns $Y_1$ need to form an invertible frame.

The initial comparison is obtained in a weak sense on transverse
cylinders. Section~\ref{sec:transport-matching} proves uniqueness for
supported planar transport equations and an estimate on fixed supported
Sobolev spaces. These results assemble the local comparisons into a
unique smooth field $G(x,[\theta])\in\operatorname{End}(\Ecal_\theta)$
on $\R^3\times\mathcal C$, where
$\mathcal C=\{[\theta]\in\C\mathrm P^2:\theta\cdot\theta=0\}$ is the
projective null conic. For a fixed ball $\mathcal B\subset\R^3$
independent of $[\theta]$, it satisfies
\[
 D_\theta G=M_{2,\theta}G-GM_{1,\theta},
 \qquad G(x,[\theta])=\Id\quad(x\notin\mathcal B).
\]
Differentiating this equation in a conjugate conic coordinate while
holding $x$ fixed gives a supported homogeneous transport equation.
Its uniqueness makes $G(x,\cdot)$ holomorphic.

Finally, the planes $H_\theta$ form a holomorphic bundle $H$ on
$\mathcal C$. It has no nonzero global holomorphic sections, while its
quotient by the null line subbundle with fiber $\C\theta$ is trivial.
These facts and the universal upper transport equation force the
displacement block of $G$ to be
scalar. The remaining upper block equation uses $\dim H_\theta=2$ to
make that scalar constant in $x$ and remove the remaining off-diagonal
block; the exterior value then gives $G=\Id$.
Equality of the transports first determines
$(\lambda+\mu)/(\lambda+2\mu)$. The remaining equations for
$\eta=\log\mu_1-\log\mu_2$ have the form
$\nabla^2\eta=B(x)\nabla\eta+q(x)I$, where
$B(x):\R^3\to\R^{3\times3}$ and the scalar $q(x)$ are smooth.
One differentiation closes a homogeneous first-order system for
$(\nabla\eta,q)$, whose zero exterior values propagate by uniqueness
for ordinary differential equations along paths.
Section~\ref{sec:coefficient-recovery} gives this rigidity and recovers
both coefficients.
